% Root proof, independently checked by /root/cycles and /root/real_weil.
\paragraph{C.3v (Refuted already in dimension two): finiteness of the ordinary divisor quotient.}
For every $N\ge2$, there is an explicitly specified excellent
normal strictly henselian local $\mathbb C$-algebra $R_N$ of
dimension $N$, with regular puncture $X_N$, for which
\begin{equation}\label{eq:C3v-ordinary-quotient}
 B^1(X_N):=\frac{\operatorname{Pic}(X_N)}
                   {\operatorname{Pic}(X_N)_{\mathrm{alg}}}
 \simeq
 \begin{cases}
  (\mathbb C,+)\oplus\mathbb Z/4,&N=2,\\
  (\mathbb C,+)^{N-1}\oplus\mathbb Z,&N\ge3.
 \end{cases}
\end{equation}
Here algebraic equivalence is ordinary algebraic equivalence
through line bundles on the fixed product with a smooth
connected curve. In particular the rational quotient is
nonzero and infinite-dimensional over $\mathbb Q$, already
in dimension two. Thus a finiteness assertion for the
degree-one quotient on 14:107 is false on this reading.
The quotient using the connected local Picard functor is a
different relation and is not refuted by this calculation.

\begin{proof}
Let
\[
 C=\{z_0^4+z_1^4+z_2^4=0\}\subset\mathbb P^2_{\mathbb C},
 \qquad m=N-2,\qquad V=\mathbb P^m\times C,
 \qquad L=\mathcal O_{\mathbb P^m}(1)\boxtimes K_C.
\]
For $m=0$, the first factor and its line bundle mean a
point and its trivial line bundle. The partial derivatives
$4z_i^3$ have no common projective zero. Adjunction gives
$K_C=\mathcal O_C(1)$, $\deg K_C=4$ and $g(C)=3$.
The canonical embedding of this plane quartic and the
standard projective embedding are projectively normal.
Their Segre section ring
\[
 S=\bigoplus_{a\ge0}H^0(V,L^a),\qquad
 R_N=(S_{S_+})^{\mathrm{sh}}
\]
is therefore a normal domain generated in degree one.
We use the strict henselization with its specified residue
field $\mathbb C$; it equals the henselization. Excellence,
normality and dimension $\dim R_N=m+2=N$ follow from
this finite-type construction. The cone puncture is the
complement of the zero section in $\operatorname{Tot}(L^{-1})$,
and is smooth. Its strict-local base change $X_N$ is regular.

Let $Z\to\operatorname{Spec}R_N$ be the base change of the
blow-up of the vertex. It is regular and proper, with sole
exceptional divisor $V$, projection $\pi:Z\to V$, and
\[
 I_V=\mathfrak m_{R_N}\mathcal O_Z,
 \qquad \mathcal O_Z(V)=\pi^*L^{-1}.
\]
Degree-one generation proves the first identity. Completion
along $V$ is the formal completion of the line bundle along
its zero section. On its $n$th infinitesimal neighborhood,
the projection splits the units onto $\mathcal O_V^*$,
while the characteristic-zero logarithm identifies the
remaining units with the additive nilpotent ideal
\[
 I_V/I_V^n=\bigoplus_{a=1}^{n-1}L^a.
\]
The logarithm is a finite polynomial at each level, so
these are actual compatible isomorphisms of sheaves of
abelian groups, not merely tangent-space identifications.
Consequently
\[
 \operatorname{Pic}(V_n)
 \simeq\operatorname{Pic}(V)\oplus
                   \bigoplus_{a=1}^{n-1}H^1(V,L^a).
\]
For $a>0$, the K\"unneth formula and Serre duality on $C$
give
\begin{align*}
 H^1(V,L)&=H^0(\mathbb P^m,\mathcal O(1))
                         \otimes H^1(C,K_C)
                  \simeq\mathbb C^{m+1},\\
 H^1(V,L^a)&=0\qquad(a\ge2).
\end{align*}
Indeed $H^1(C,K_C^a)$ is dual to
$H^0(C,K_C^{1-a})$, whose degree is negative for $a\ge2$;
the higher cohomology of positive projective-space line
bundles contributes nothing in degree one.

There is no inverse-limit ambiguity here. Global units
on $V_n$ split as $\mathbb C^*$ times the finite direct
sum of the $H^0(V,L^a)$ under the logarithm. Their
transition maps are surjective, so the corresponding
$\varprojlim^1$ is zero. Equivalently, isomorphisms of
successive line bundles can be chosen compatibly by
these surjections. Thus the formal Picard group is
\[
 \operatorname{Pic}(\widehat Z)
       \simeq\operatorname{Pic}(V)\oplus\mathbb C^{m+1}.
\]
Grothendieck existence, in the precise form of
\citet[Theorem 30.27.1, Tag 088E]{rootStacksExistence},
algebraizes every such compatible line bundle to the
proper scheme $Z_{\widehat R_N}$ and is fully faithful.
The algebraized coherent sheaf is invertible: it is
invertible near the closed fibre by Nakayama's lemma,
and the image of its closed nonfree locus would be a
nonempty closed subset of the local base missing its
closed point, which properness excludes.

The excellent regular base changes used here keep
$Z_{\widehat R_N}$ and its puncture regular. Cartier
divisor localization therefore identifies
\[
 \operatorname{Pic}(X_{\widehat R_N})
   \simeq
   \bigl(\operatorname{Pic}(V)/\mathbb Z[L]\bigr)
                                      \oplus\mathbb C^{m+1}.
\]
The exceptional divisor contributes exactly $[L^{-1}]$
and has zero additive coordinate. Finally,
\citet[Section III, Theorem 3, p.~569, and its unnumbered
Picard corollary, p.~573]{D_Elkik1973}
gives a bijection
$\operatorname{Pic}(X_N)\to\operatorname{Pic}(X_{\widehat R_N})$.
Its hypotheses are the noetherian henselian pair
$(R_N,\mathfrak m_{R_N})$ and its completion; there is no
dimension-at-least-three assumption. We have proved the
same direct-sum formula for the actual punctured
henselian scheme $X_N$.

We now identify ordinary algebraic equivalence without
replacing it by formal-family equivalence. Suppose
$\mathcal M$ is an actual line bundle on $X_N\times T$,
where $T$ is a smooth connected curve. Compactify $T$
to a smooth proper connected curve $\overline T$.
The schemes $X_N\times\overline T$ and
$Z\times\overline T$ are regular and noetherian.
Closing a Cartier divisor first across
$X_N\times(\overline T\setminus T)$ and then across
$V\times\overline T$ extends $\mathcal M$ to a line
bundle on $Z\times\overline T$. Restrict it to
$V_2\times\overline T$. The split square-zero
description above is functorial in a parameter scheme;
it gives the Picard group scheme decomposition
\[
 \mathbf{Pic}_{V_2/\mathbb C}
   \simeq \mathbf{Pic}_{V/\mathbb C}
                       \times\mathbb G_a^{m+1}.
\]
The resulting map from the proper connected curve
$\overline T$ to the affine vector group is constant.
Its two fibre differences therefore have zero additive
coordinate. Their $\operatorname{Pic}(V)$ coordinates
differ by an element of $\operatorname{Pic}^0(V)(\mathbb C)$.
The ambiguity of extension by exceptional divisors has
zero additive coordinate and vanishes modulo $\mathbb Z[L]$.
It follows, for every actual family and hence every sum
of family differences, that
\[
 \operatorname{Pic}(X_N)_{\mathrm{alg}}
       \subseteq\rho^*\operatorname{Pic}^0(V)(\mathbb C),
\]
where $\rho:X_N\to V$ is the radial projection.
Conversely a Poincar\'e line bundle on
$V\times\operatorname{Pic}^0(V)$, pulled back by $\rho$,
gives these ordinary algebraic classes. Any two points
of the abelian parameter can be joined by a smooth
connected curve (or a chain of such curves). Also
$\mathbb Z[L]\cap\operatorname{Pic}^0(V)=0$, since a
nonzero power of an ample line bundle is not numerically
trivial. Thus this subgroup is an actual embedded copy
of $\operatorname{Pic}^0(C)(\mathbb C)$.

Taking the quotient gives
\[
 B^1(X_N)\simeq
       \bigl(\operatorname{NS}(V)/\mathbb Z[L]\bigr)
                                      \oplus\mathbb C^{m+1}.
\]
If $m=0$, $\operatorname{NS}(C)=\mathbb Z$ by degree and
$[K_C]$ has degree four. If $m\ge1$, the projective
bundle formula gives $\operatorname{NS}(V)=\mathbb Z^2$,
and $[L]=(1,4)$ is primitive, leaving one free copy
of $\mathbb Z$. This proves \eqref{eq:C3v-ordinary-quotient}.

For completeness, we specify one non-torsion divisor
class in the smallest counterexample, rather than only
its ambient group. Choose $\eta\in\mathbb C$ with
$\eta^4=-1$ and let $c_0=(1:\eta:0)\in C$.
The function $t=z_2/z_0$ is a uniformizer at $c_0$,
since the derivative of $1+(z_1/z_0)^4+t^4$ with respect
to $z_1/z_0$ is $4\eta^3\ne0$ there. Take
$U=C\setminus\{c_0\}$ and an affine neighborhood
$W\subset\{z_0\ne0\}$ of $c_0$ excluding the other
three zeros of $t$. The regular differential
$dt/t$ on $U\cap W$ represents a nonzero class
$\alpha\in H^1(C,K_C)$: its Serre-duality residue
pairing with $1\in H^0(C,\mathcal O_C)$ is one.
On the formal total space of $K_C^{-1}$, contract
$dt/t$ with the tautological radial vector to obtain
$a\in I_V(U\cap W)$. Glue the trivial formal line
bundles by the transition $\exp(a)$; this is a
compatible formal unit, and its additive Picard
coordinate is exactly $\alpha$. Proper formal existence
algebraizes it, and Elkik's bijection specifies a unique
line-bundle class $M$ on $X_2$. This construction fixes
its class independently of all auxiliary algebraization
choices. Its image in $B^1(X_2)$ has additive coordinate
$\alpha\ne0$, so every nonzero integer multiple remains
nonzero. A Cartier divisor representing $M$ is the
required ordinary cycle counterexample.
\end{proof}

The finite certificate
\texttt{code/check\_ordinary\_divisor\_quotient.py}
checks the smooth quartic equations, the chosen local
parameter, the residue coefficient, the indicated
cohomology dimensions in its stated finite input range,
and the two N\'eron--Severi lattice quotients. The
formal descent, the all-dimension formula and the
non-torsion conclusion follow from the proof, not
from extrapolation of that computation.
